\subsection{Elastic F-table boundary switch}
\label{sec:verification:elastic-f-table}

This test isolates prescribed F-table deformation, boundary-condition-table switching, reaction output, and box-average stress output.  A unit elastic cube is stretched in the $x$ direction.  During the first segment all faces are kinematically constrained; after $t=0.5$ the lateral faces are released while the $x$ stretch continues.

\begin{figure}[h]
  \centering
  \includegraphics[width=0.72\linewidth]{\CaseSourceDir/setup_elastic_f_table.png}
  \caption{Elastic F-table setup.  The first interval uses constrained lateral faces, then the boundary-condition table releases the lateral faces while the axial F-table remains active.}
\end{figure}

For small strain, the expected axial stress is piecewise linear.  Before lateral release,
\begin{equation}
  \sigma_{xx}=M\epsilon_{xx}, \qquad
  M=\frac{E(1-\nu)}{(1+\nu)(1-2\nu)}.
\end{equation}
After release, incremental stress changes are compared with the Young's modulus $E$ while preserving stress continuity at the switching time.  The metric is RMS and maximum axial box-average stress error.

\IfFileExists{\CaseOutputDir/elastic_f_table_results.tex}{\input{\CaseOutputDir/elastic_f_table_results.tex}}{\paragraph{Generated results.} Results are produced by \texttt{postProcess\_elasticBlockUni.py} after the GEOS job finishes.}
